% Questions for sections/01_intro.tex, 01b_attacks.tex, 02_classical_caesar.tex (end-sem rewrite).
% Numerical answers checked by script (scratchpad/r1/verify.py). Long practice set: questions/99_practice.tex, part A.
\section{Security Concepts and the Caesar Cipher}\label{q:01_intro}
\noindent{\small Notes: \S\ref{sec:netmodel}--\S\ref{sec:schemes}, \S\ref{sec:attacks}, \S\ref{sec:classical}--\S\ref{sec:bruteforce}.}

\subsection*{Check your understanding}
\begin{enumerate}[leftmargin=1.6em,itemsep=3pt]
  \item \textbf{In the model for network security, what are the two inputs of the transformation $T$, and why is it called $T$ rather than $E$?} $T$ takes the message and secret information (a key) and outputs a secure message; the receiver applies the complementary transformation with its own secret information. It is called $T$ because encryption is only one possible security-related transformation; steganography, which hides the message inside a carrier, is another.
  \item \textbf{List the four tasks in designing a security service and link each to the model.} (1) Design an algorithm for the transformation $T$; (2) generate the secret information used with it; (3) develop methods to distribute and share the secret information (key distribution, often via the trusted third party); (4) specify a protocol by which the principals use the algorithm and the secret information to achieve a particular service. The first three come from the components $T$, secret info and TTP; the fourth fixes how they are used together.
  \item \textbf{Differentiate the model for network security from the network access security model.} The first protects a message in transit between two principals against an opponent on the channel, using a transformation with secret information. The second protects an information system against opponents (hackers, viruses, worms) who try to get in through an access channel, using a gatekeeper function (firewall, login) and internal security controls (antivirus, intrusion detection).
  \item \textbf{Name the five ingredients of the symmetric cipher model and the two requirements for its secure use.} Plaintext, encryption algorithm, secret key, ciphertext and decryption algorithm. The algorithm must be strong enough that an opponent who knows it and holds ciphertexts (even with their plaintexts) cannot find the key or other plaintexts, and the key must reach sender and receiver securely and be kept secret.
  \item \textbf{Define unconditionally secure and computationally secure, with an example of each.} Unconditionally secure: the ciphertext does not contain enough information to determine the plaintext uniquely, however much ciphertext is available, even against unlimited computing power (one-time pad). Computationally secure: the cost of breaking exceeds the value of the information, or the time to break exceeds its useful lifetime (AES, RSA).
  \item \textbf{How long does an average exhaustive search of a 56-bit key take at $10^9$ keys per second, and why is a 128-bit key considered safe?} On average half the keys are tried: $2^{55}/10^9 \approx 3.6 \times 10^7$ s $\approx$ \textbf{1.14 years}, which is within reach (and an hour at $10^{13}$/s). A 128-bit key needs $2^{127}/10^9 \approx 5.4 \times 10^{21}$ years, far longer than the useful life of any information.
  \item \textbf{Differentiate passive and active attacks. Why are passive attacks hard to detect but preventable?} A passive attack only obtains information (snooping, traffic analysis) and threatens confidentiality; an active attack changes data or harms the system (modification, masquerading, replaying, repudiation, DoS) and threatens integrity or availability. Since a passive attack changes nothing, there is nothing to detect, but encipherment makes the captured data useless; active attacks leave visible effects and are easier to detect than to prevent.
  \item \textbf{Why is repudiation different from all other attacks, and what counters it?} It is carried out by one of the two communicating parties (the sender denies sending, or the receiver denies receiving), not by an outsider. Encryption cannot help, since both parties hold the key; it is countered by the nonrepudiation service, implemented by digital signatures, data integrity and notarization.
  \item \textbf{Which X.800 service and mechanism counter (a) replaying, (b) masquerading, (c) traffic analysis?} (a) Data integrity in its anti-replay form, using sequence numbers, timestamps or nonces with an integrity check; (b) authentication, by encipherment, digital signature or authentication exchange; (c) traffic-flow confidentiality, by traffic padding and routing control.
  \item \textbf{Differentiate a security service from a security mechanism, and explain why the digital signature appears under three services.} A service is what is provided (confidentiality, integrity, authentication, nonrepudiation, access control); a mechanism is how it is provided (encipherment, signature, notarization, \dots). A signature binds the data to the holder of one private key, so a change breaks it (integrity), a valid signature identifies the signer (authentication), and the signer cannot later deny it (nonrepudiation).
  \item \textbf{State Kerckhoff's principle and give two reasons for it.} Assume the adversary knows the encryption and decryption algorithms; security must rest only on the secrecy of the key. Algorithms leak through products and people and can be reverse-engineered, whereas a compromised key is cheap to replace; and published algorithms are tested by the whole research community.
  \item \textbf{Order the four cryptanalysis attacks by the information Eve has. Which is most likely, and why?} Ciphertext-only $<$ known-plaintext $<$ chosen-plaintext / chosen-ciphertext. Ciphertext-only is the most likely, because intercepting ciphertext is all it needs; the chosen attacks need access to Alice's encryptor or Bob's decryptor and are the least likely but the most powerful.
  \item \textbf{Given only a long ciphertext, how can you tell whether a substitution or a transposition cipher was used?} Count the letter frequencies. A transposition only rearranges letters, so the ciphertext has the normal English frequency profile (e, t, a most common) without readable words; a monoalphabetic substitution shows the same profile moved onto different letters.
  \item \textbf{Encrypt ``attack'' with a shift cipher, $k = 11$, and check by decrypting.} $p = 0, 19, 19, 0, 2, 10$; $p + 11 = 11, 30, 30, 11, 13, 21$; mod 26: $11, 4, 4, 11, 13, 21$ = \textbf{\ttx{LEELNV}}. Decrypting, $\ttx{E} = 4$ gives $4 - 11 = -7 \equiv 19 = \ttx{t}$, and the whole word returns to \ttx{attack} \checkmark.
  \item \textbf{Why does the shift cipher have 25 useful keys, and how many trials does a brute-force attack need on average? In a long ciphertext the most frequent letter is \ttx{K}; which key do you try first?} Keys are only defined modulo 26 ($k$ and $k + 26$ act alike), and $k = 0$ leaves the text unchanged, leaving $1, \dots, 25$; with equally likely keys the expected number of trials is $(1 + \cdots + 25)/25 = 13$. Taking \ttx{K} $= 10$ as \ttx{e} $= 4$ gives the first key to try, $k = (10 - 4) \bmod 26 = \mathbf{6}$.
\end{enumerate}
